\appendix

\section{Implementation Details}
\label{sec:appendix:implementation-details}

During evaluation, speculation is run at each realized tool-call boundary of the
agent trajectory. The speculator receives the prefix before the tool call and
predicts a candidate call, but the environment executes only the agent's actual
call. This keeps the evaluated trajectory identical to the ordinary agent
trajectory and makes each speculation label the current agent's own next action.
During RL, we use the same prefix--next-call construction to derive speculation
examples from freshly sampled agent rollouts, as described in
Section~\ref{sec:method:dual-mode-rollout}.

\subsection{Training Hyperparameters}
\label{sec:appendix:training-hyperparameters}

Table~\ref{tab:training-hyperparameters} lists the default training settings
used for joint agent-speculator RL. Unless stated otherwise, ablations use the
same rollout budget, learning rate, schedule, and evaluation protocol as the
corresponding full run.

\begin{table}[h]
\centering
\setlength{\tabcolsep}{4pt}
\resizebox{\columnwidth}{!}{%
\begin{tabular}{@{}ll@{}}
\toprule
Setting & Value \\
\midrule
Backbone & Qwen3-4B, Qwen3.5-4B \\
RL method & DAPO-style GPO \\
$G_{\mathrm{ag}}$ & 8 \\
$G_{\mathrm{sp}}$ & 8 \\
Schedule & 4 agent / 8 speculator \\
Reset & Each mode switch \\
Learning rate & $3\times10^{-6}$ \\
Training length limit & 16K tokens (SearchQA); 32K tokens ($\tau$-bench) \\
RL budget & 200 steps per setting \\
Training hardware & 8 H100 GPUs (SearchQA); 8 H200 GPUs (ToolScale) \\
SFT data & Successful trajectories and speculation examples \\
SFT target & Agent trajectories; next calls under the speculation suffix \\
\bottomrule
\end{tabular}
}
\caption{Default training hyperparameters.}
\label{tab:training-hyperparameters}
\end{table}

\subsection{Prompt and Tool Schema Details}
\label{sec:appendix:prompt-and-tool-schema-details}

The fixed speculation suffix used in all experiments is:
\begin{quote}
\small\ttfamily
\textless think\textgreater{} Okay, let's see. The user provided what I need.
I'll look it up. The next step is to make the tool call.
\textless/think\textgreater{}
\end{quote}

For each speculation query, the suffix is appended after the intermediate
trajectory prefix and before the model decodes the candidate tool call. We use
the same suffix for off-the-shelf speculator evaluation, SFT warmup example
construction, joint agent-speculator RL, and final speculation evaluation. The
suffix does not name a benchmark, tool, argument key, or target answer.

We represent every tool call as a structured action $a=(n,\sigma)$, where $n$
is the tool name and $\sigma$ is a dictionary of named arguments. SearchQA
actions contain a search tool name and query string. $\tau$-bench-style actions
use the API name and schema-constrained argument dictionary supplied by the
environment. Calls that cannot be parsed into the expected structured format are
treated as invalid predictions and receive zero speculation reward.

For Hit@1 evaluation, a prediction is correct only if the parsed tool name and
the complete parsed argument dictionary exactly match the agent's eventual call.
For the shaped speculation reward used during RL, we first require the tool name
to match and then compute macro token-F1 over target argument keys, as defined
in Equation~\ref{eq:reward}. Missing predicted argument keys receive zero score,
and extra predicted keys do not compensate for missing target keys. This keeps
the training reward dense enough to provide partial credit for close arguments,
while the deployment reuse metric remains exact structured-call equality.

\section{Artifact Licenses}
\label{sec:appendix:artifact-licenses}

Table~\ref{tab:artifact-licenses} summarizes the licenses reported by the
upstream publishers for the external artifacts used in our experiments. We use
these artifacts for research evaluation and model training under their published
terms, and we retain the corresponding citations and license notices when
redistributing metadata or scripts.

\begin{table}[h]
\centering
\small
\setlength{\tabcolsep}{4pt}
\begin{tabular}{@{}p{0.62\columnwidth}p{0.30\columnwidth}@{}}
\toprule
Artifact & Reported license \\
\midrule
Qwen3-0.6B, Qwen3-1.7B, Qwen3-4B &
Apache-2.0 \\
Qwen3.5-0.8B, Qwen3.5-2B, Qwen3.5-4B &
Apache-2.0 \\
Qwen3-Embedding-0.6B &
Apache-2.0 \\
HotpotQA &
CC BY-SA 4.0 \\
MuSiQue &
CC BY 4.0 \\
BrowseComp-Plus &
MIT \\
$\tau$-bench / $\tau^2$-bench &
MIT \\
ToolScale &
Upstream dataset terms \\
\bottomrule
\end{tabular}
\caption{Licenses or access terms reported by the upstream model, dataset, or
benchmark publishers for artifacts used in the experiments.}
\label{tab:artifact-licenses}
\end{table}

\section{Data Statistics}
\label{sec:appendix:data-statistics}

Table~\ref{tab:data-statistics} reports the train, evaluation, and SFT warmup
counts used in our experiments. SFT counts are separated into agent-trajectory
rows and speculation rows; slash-separated counts denote Qwen3-4B / Qwen3.5-4B
when both warmup files are used. Following the subset-based evaluation practice
in R1-Searcher~\citep{song2025r1searcher}, SearchQA evaluation uses
500 examples per dataset sampled from the FlashRAG development splits, rather
than the full public development or test sets. We do not assume that these are
the identical examples released by R1-Searcher unless their instance IDs match
explicitly in metadata.

\begin{table}[h]
\centering
\small
\setlength{\tabcolsep}{4pt}
\resizebox{\columnwidth}{!}{%
\begin{tabular}{@{}lcccc@{}}
\toprule
Resource & Train & Eval & SFT agent & SFT spec \\
\midrule
FlashQA / SearchQA total & 283,338 & 2,500 subset & 1,658 / 1,351 & 1,658 / 1,351 \\
HotpotQA & -- & 500 subset & -- & -- \\
MuSiQue & -- & 500 subset & -- & -- \\
BrowseComp-Plus & 680 & 150 & -- & -- \\
ToolScale & 4,063 & -- & 0 & 5,000 \\
$\tau$-bench Airline & -- & 50 & -- & -- \\
$\tau$-bench Retail & -- & 114 & -- & -- \\
\bottomrule
\end{tabular}
}
\caption{Dataset, benchmark, and SFT warmup counts. The ToolScale warmup source
contains speculation prompt--label rows only.}
\label{tab:data-statistics}
\end{table}

\section{LLM Usage Statement}
\label{sec:appendix:llm-usage}

We used large language models to assist with paper grammar editing and with code
implementation and debugging. 